% --- Archon physics formalization source begin ---
% archon:physics
% archon:chemistry
% archon:covers IChO2026Problems/problem_icho_2026_t5_a3.lean
% archon:source-report reports/icho_2026/problem_icho_2026_t5_a3.source.json
% archon:problem-id icho_2026_t5
% archon:part-id T5-A3
% archon:previous-part icho_2026_t5_a1
% archon:previous-part-policy natural_language_prerequisite_only
% archon:previous-part icho_2026_t5_a2
% archon:previous-part-policy natural_language_prerequisite_only

\chapter{Icho Chemistry problem icho\_2026\_t5\_a3}
\label{ch:IChO2026Problems_problem_icho_2026_t5_a3}

\paragraph{Problem source.}
\#\# Source
58th International Chemistry Olympiad, Tashkent, Uzbekistan, 2026.
Official-materials index: https://www.icho2026.uz/problems
English problem PDF: ../icho\_2026\_source/raw/theory\_problem.pdf
English solution/rubric PDF: ../icho\_2026\_source/raw/theory\_solution.pdf
Problem PDF page 46; printed local question page 3; solution starts on PDF page 42.
The page PNG is primary visual evidence for formulas, structures, tables, and figures.

\#\# T5. Cardiolipins
T5. Cardiolipins 6\% In this problem you are asked to be creative and build molecules from fragments. For instance, if you are required to draw a chiral phospholipid such as W using fragments a-e, it can be assembled as shown below. Cardiolipins are a family of acyclic phospholipids differing in fatty acid residues. They are essential for energy production and membrane stability. Defects can lead to certain diseases. PL1 belongs to this family. Using only the structural elements a-d in the quantities stated below, it is possible to assemble the non-ionised form of PL1. R is a hydrocarbon substituent in the fatty acid structure. PL1 does not contain any peroxide bonds. 5.1 Tick one correct statement regarding the value of n (number of type a fragments). (a) n is an even number, (b) n is an odd number, (c) n can be either an odd or an even number. Some cardiolipins are chiral molecules, even if they contain four identical fatty acid residues, as in the case of PL1. One enantiomer of PL1 is found in prokaryotes and eukaryotes, and the other only in archaea. All other diastereomers of PL1 are achiral molecules, since they have a plane of symmetry. Do not consider the chirality of phosphorus atoms as stereocentres. PL1 is a diprotic acid with the same acidic groups. Its second deprotonation is less favourable (𝑝𝐾𝑎2 ≫𝑝𝐾𝑎1) due to the special, stable structure of monoanion Y. 5.2 Draw the structure of one enantiomer of PL1. Draw the structure of Y in a way that shows the stabilisation that explains the difference between 𝑝𝐾𝑎1 and 𝑝𝐾𝑎2. Use the abbreviation R for the fatty acid residues. A series of experiments was conducted to identify the fatty acid (RCOOH) in PL1 found predominantly in mammalian hearts. During reductive ozonolysis, RCOOH forms three different organic products in equimolar amounts.

\#\# Current subquestion 5.3
Determine the molecular formula of the fatty acid (RCOOH), if the non-ionised form of PL1 contains 255 𝜎and 𝜋bonds between atoms in total. Support your answer with calculations. If you were unable to find the structural formula of PL1, you can use a-d fragments from 5.1.

\#\# Dataset metadata
IChO 2026; T5; raw rubric points=4; kind=theory; formalization\_ready=true.

\paragraph{Shared problem context.}
T5. Cardiolipins 6\% In this problem you are asked to be creative and build molecules from fragments. For instance, if you are required to draw a chiral phospholipid such as W using fragments a-e, it can be assembled as shown below. Cardiolipins are a family of acyclic phospholipids differing in fatty acid residues. They are essential for energy production and membrane stability. Defects can lead to certain diseases. PL1 belongs to this family. Using only the structural elements a-d in the quantities stated below, it is possible to assemble the non-ionised form of PL1. R is a hydrocarbon substituent in the fatty acid structure. PL1 does not contain any peroxide bonds. 5.1 Tick one correct statement regarding the value of n (number of type a fragments). (a) n is an even number, (b) n is an odd number, (c) n can be either an odd or an even number. Some cardiolipins are chiral molecules, even if they contain four identical fatty acid residues, as in the case of PL1. One enantiomer of PL1 is found in prokaryotes and eukaryotes, and the other only in archaea. All other diastereomers of PL1 are achiral molecules, since they have a plane of symmetry. Do not consider the chirality of phosphorus atoms as stereocentres. PL1 is a diprotic acid with the same acidic groups. Its second deprotonation is less favourable (𝑝𝐾𝑎2 ≫𝑝𝐾𝑎1) due to the special, stable structure of monoanion Y. 5.2 Draw the structure of one enantiomer of PL1. Draw the structure of Y in a way that shows the stabilisation that explains the difference between 𝑝𝐾𝑎1 and 𝑝𝐾𝑎2. Use the abbreviation R for the fatty acid residues. A series of experiments was conducted to identify the fatty acid (RCOOH) in PL1 found predominantly in mammalian hearts. During reductive ozonolysis, RCOOH forms three different organic products in equimolar amounts.

\paragraph{Current subquestion.}
Determine the molecular formula of the fatty acid (RCOOH), if the non-ionised form of PL1 contains 255 𝜎and 𝜋bonds between atoms in total. Support your answer with calculations. If you were unable to find the structural formula of PL1, you can use a-d fragments from 5.1.

\paragraph{Recorded answer/context.}
Correct scheme of reductive ozonolysis of RCOOH is as follows: So, 𝑁𝐶=𝐶= 2. If we take R as 𝐶𝑥𝐻𝑦, then: y=(2x+1)-𝑁𝐶=𝐶∙2=2x-3.(eq.3.1) Then the formula for calculating the number of bonds in PL1 can be written as follows: 0.5∙1+5∙2+11.5∙3+4∙(3+2x+0.5y)=255, 4x+y=99.(eq.3.2) Combining eq.3.1 and eq.3.2: 4x+y=4x+2x-3=6x-3=99 ⟹x=17, y=31. RCOOH −C17H31COOH or C18H32O2 For correct eq.3.2: 2 points. For correct number of double bonds: 1 point. For correct formula of C17H31COOH: 1 point. 4.0 pt

\paragraph{Figure/image paths.}
\begin{itemize}
\item ../icho\_2026\_source/image/T5\_page-3.png
\item ../icho\_2026\_source/image/T5\_page-2.png
\item ../icho\_2026\_source/image/T5\_page-1.png
\end{itemize}

\paragraph{Reusable previous-part conclusions.}
\begin{itemize}
\item Source T5-A1. Question: Tick one correct statement regarding the value of n (number of type a fragments). (a) n is an even number, (b) n is an odd number, (c) n can be either an odd or an even number. Reusable conclusions: Each bond involves two atoms, so the total number of atoms with broken bonds (indicated by the wavy line on the fragments of the molecule) must be an even number. For fragments with a known number of repeats in the PL1 molecule, the total number of atoms with broken bonds is 17 (2·2 + 3·3 + 4). Thus the total number of hydrogen atoms with broken bonds must be an odd number. Therefore, n is an odd number. Correct answer: 1 point. 1.0 pt Policy: natural\_language\_prerequisite\_only; the current target and later answers must not be assumed
\item Source T5-A2. Question: Draw the structure of one enantiomer of PL1. Draw the structure of Y in a way that shows the stabilisation that explains the difference between 𝑝𝐾𝑎1 and 𝑝𝐾𝑎2. Use the abbreviation R for the fatty acid residues. Reusable conclusions: Visual-answer notice: the official solution contains essential ticks, structures, tables, graphs, or equations that PDF plain-text extraction may omit. Consult the cited solution PDF page.

To obtain the diastereomer of PL1 with a plane of symmetry, we must construct a symmetric structure from the given fragments. Glycerol backbones can only be linked through phosphate groups since there is no peroxide bond present. The number of glycerol moieties suggests that one glycerol occupies the central position with the remaining fragments arranged around it to achieve an overall plane of symmetry. The number of a fragments in PL1 is n = 3·3 – 2·2 – 4 = 1. So, PL1 has 𝑑2𝑐−𝑏−𝑐(𝑎) −𝑏−𝑐𝑑2 skeleton. Two options, presented below, can have a plane of symmetry: The PL0 structure is achiral, so it cannot be PL1. All possible stereoisomers of PL1 are as follows: SOLUTION: Enantiomer 2 is found in prokaryotes and eukaryotes. These organisms use glycerol 3-phosphate for cardiolipin synthesis, while archaea use glycerol 1phosphate for this purposes and synthesise enantiomer 1: 6.0 pt SOLUTION: PL1 For drawing of enantiomer 1 or 2 as PL1: 4 points. For drawing the correct structure of PL1 with wrong or missing stereochemistry: 3 points. For drawing achiral PL0 structure: 2 points. Any other structure that uses exactly 1×a, 2×b, 3×c, 4×d and contains no peroxide bond: 1 point. Intramolecular hydrogen bond between the phosphoric acid residue and the alcohol group significantly reduces the 𝑝𝐾𝑎2 value. For any correct tautomer of Y: 2 points. For a tautomer of Y (not shown above) with a hydrogen bond between the two phosphate groups: 1 point Correctly drawn structures with incorrect PL1 structure are also accepted. For structures containing hydrogen bond between phosphoric acid residues and esters 0 point, because they are possible for each phosphoric acid residue even in molecular form and don’t make difference between 𝑝𝐾𝑎1 and 𝑝𝐾𝑎2. Policy: natural\_language\_prerequisite\_only; the current target and later answers must not be assumed
\end{itemize}

\paragraph{Formalization target.}
The verified Lean formalization is in `IChO2026Problems/problem\_icho\_2026\_t5\_a3.lean`,
with its theorem and lemma proofs complete.
The Lean declarations must preserve every mathematical object, quantifier, hypothesis, side condition, approximation/error guarantee, and requested conclusion from the source statement.
The implementation reuses the pinned Mathlib, Physlib, and Chemistry libraries,
with local definitions only for source-specific objects.

\begin{theorem}[Icho Chemistry formalization target]
\label{thm:physics:icho_2026_t5_a3:target}
The mapped declarations encode the stated calculation and its verified conclusion.
\end{theorem}
\begin{proof}
The proof unfolds the source-specific definitions and establishes the mapped
conclusion from the encoded hypotheses.
\end{proof}

\paragraph{Informal proof (official marking scheme).}
Writing the hydrocarbon substituent as $R = C_xH_y$: since $R$ is an
acyclic monovalent hydrocarbon group whose only unsaturation is
$N_{C=C}$ double bonds, the saturated formula $C_xH_{2x+1}$ loses two
hydrogens per double bond,
\begin{equation}
y = (2x+1) - 2N_{C=C}. \tag{eq.\,3.1}
\end{equation}
Reductive ozonolysis cleaves each C=C once, so an acyclic acid with
$N_{C=C}$ double bonds affords $N_{C=C}+1$ carbonyl products, one
molecule of each per molecule of acid (hence automatically equimolar);
the observed three products force $N_{C=C} + 1 = 3$, i.e.\
$N_{C=C} = 2$ and $y = 2x - 3$.

Counting $\sigma$- and $\pi$-bonds of non-ionised PL1 fragmentwise
(each fragment contributes its internal bonds, a double bond counting
as $2$, plus $\tfrac12$ per open valence, so that each inter-fragment
bond is counted once in total) with $1\times a$, $2\times b$,
$3\times c$, $4\times d$ gives
\[
0.5\cdot 1 + 5\cdot 2 + 11.5\cdot 3 + 4\cdot(3 + 2x + 0.5y) = 255,
\]
which simplifies to
\begin{equation}
4x + y = 99. \tag{eq.\,3.2}
\end{equation}
Combining (eq.\,3.1) and (eq.\,3.2): $4x + (2x - 3) = 6x - 3 = 99$,
hence $x = 17$ and $y = 31$. The fatty acid is $C_{17}H_{31}COOH$,
molecular formula $C_{18}H_{32}O_2$.

\begin{lemma}[Ozonolysis forces two double bonds]
\label{thm:physics:icho_2026_t5_a3:ozonolysis_gives_two_double_bonds}
\lean{IChO2026T5A3.ozonolysis_gives_two_double_bonds}
\leanok
If the reductive ozonolysis of the acyclic fatty acid produces
\textup{three} organic products, then the acid contains exactly two
C=C double bonds: \texttt{ozonolysisProductCount} $d = d + 1 = 3$
forces $d = 2$.
\end{lemma}
\begin{proof}
\leanok
Unfold \texttt{ozonolysisProductCount} to $d + 1 = 3$ and close by
\texttt{omega}.
\end{proof}

\begin{lemma}[Bond-count equation $4x + y = 99$]
\label{thm:physics:icho_2026_t5_a3:bond_count_equation}
\lean{IChO2026T5A3.bond_count_equation}
\leanok
The 255 $\sigma$-and-$\pi$-bond count of non-ionised PL1 yields the
marking scheme's equation~\textup{(eq.\,3.2)},
$4x + y = 99$ over $\mathbb{Q}$.
\end{lemma}
\begin{proof}
\leanok
Unfold \texttt{pl1TotalBonds} and the fragment bond counts; the
hypothesis becomes $57 + 8x + 2y = 255$ over $\mathbb{Q}$, and
\texttt{linarith} derives $4x + y = 99$.
\end{proof}

\begin{theorem}[The fatty acid is $C_{17}H_{31}COOH$]
\label{thm:physics:icho_2026_t5_a3:icho_2026_t5_a3}
\lean{IChO2026T5A3.icho_2026_t5_a3}
\leanok
For a hydrocarbon substituent $R = C_xH_y$ ($x \geq 1$) that is acyclic
and monovalent with $d$ C=C bonds as its only unsaturation
($y + 2d = 2x + 1$), whose acid affords three ozonolysis products, and
for which non-ionised PL1 has 255 $\sigma$-and-$\pi$ bonds, one has
$x = 17$, $y = 31$, and the acid's molecular formula is
$C_{18}H_{32}O_2$.
\end{theorem}
\begin{proof}
\leanok
Combine the two lemmas above: $d = 2$ in (eq.\,3.1) gives
$y = 2x - 3$; together with $4x + y = 99$ (transported to
$\mathbb{N}$ by \texttt{exact\_mod\_cast}) \texttt{omega} yields
$x = 17$ and $y = 31$; the formula check
\texttt{rcoohFormula} $17\ 31 = \langle 18, 32, 2\rangle$ is
definitional (\texttt{rfl}).
\end{proof}
% --- Archon physics formalization source end ---
